\chapter*{Annexe : Code source de la fonctionnalité de recueil de réclamations}
\addcontentsline{toc}{chapter}{Annexe : Code source de la fonctionnalité de recueil de réclamations}

%===========================================
\section*{Annexe A : Back-end Spring Boot — Gestion des réclamations}
%===========================================

\subsection*{ReclamationController.java}
\begin{lstlisting}[language=Java]
package com.bewize.monitorbackend.controller;

import com.bewize.monitorbackend.domains.ReclamationEntity;
import com.bewize.monitorbackend.dto.ReclamationDto;
import com.bewize.monitorbackend.service.ReclamationService;
import io.swagger.v3.oas.annotations.Operation;
import io.swagger.v3.oas.annotations.tags.Tag;
import lombok.RequiredArgsConstructor;
import org.springframework.http.ResponseEntity;
import org.springframework.web.bind.annotation.*;
import java.util.List;
import java.util.Map;

@Tag(name = "Reclamations")
@RestController
@RequiredArgsConstructor
@RequestMapping("/reclamations")
@CrossOrigin(originPatterns = "*")
public class ReclamationController {

    private final ReclamationService reclamationService;

    @Operation(summary = "Soumettre une nouvelle reclamation")
    @PostMapping
    public ResponseEntity<Map<String, Object>> creer(
            @RequestBody ReclamationDto dto) {
        ReclamationEntity saved = reclamationService.creer(dto);
        return ResponseEntity.ok(Map.of(
                "success", true,
                "id",      saved.getId(),
                "message", "Reclamation enregistree avec succes"
        ));
    }

    @Operation(summary = "Lister toutes les reclamations")
    @GetMapping
    public ResponseEntity<List<ReclamationEntity>> lister() {
        return ResponseEntity.ok(reclamationService.listerToutes());
    }
}
\end{lstlisting}

\newpage
\subsection*{ReclamationService.java}
\begin{lstlisting}[language=Java]
package com.bewize.monitorbackend.service;

import com.bewize.monitorbackend.domains.ReclamationEntity;
import com.bewize.monitorbackend.dto.ReclamationDto;
import com.bewize.monitorbackend.repository.ReclamationRepository;
import lombok.RequiredArgsConstructor;
import org.springframework.stereotype.Service;
import java.util.List;

@Service
@RequiredArgsConstructor
public class ReclamationService {

    private final ReclamationRepository reclamationRepository;

    public ReclamationEntity creer(ReclamationDto dto) {
        ReclamationEntity entity = ReclamationEntity.builder()
                .prenom(dto.getPrenom())
                .nom(dto.getNom())
                .email(dto.getEmail())
                .numeroEtudiant(dto.getNumeroEtudiant())
                .statut(dto.getStatut())
                .typeReclamation(dto.getTypeReclamation())
                .objet(dto.getObjet())
                .description(dto.getDescription())
                .nomFichier(dto.getNomFichier())
                .build();
        return reclamationRepository.save(entity);
    }

    public List<ReclamationEntity> listerToutes() {
        return reclamationRepository.findAll();
    }
}
\end{lstlisting}

\newpage
\subsection*{ReclamationRepository.java}
\begin{lstlisting}[language=Java]
package com.bewize.monitorbackend.repository;

import com.bewize.monitorbackend.domains.ReclamationEntity;
import org.springframework.data.jpa.repository.JpaRepository;
import org.springframework.stereotype.Repository;

@Repository
public interface ReclamationRepository
        extends JpaRepository<ReclamationEntity, Long> {
}
\end{lstlisting}

%===========================================
\newpage
\section*{Annexe B : Formulaire WebView — Saisie des réclamations}
%===========================================

\subsection*{CentreReclamations.tsx}
\begin{lstlisting}[language=Java]
import { useState, useRef } from "react";
import { insertReclamation } from "../lib/reclamations";

interface FormState {
  prenom: string;
  nom: string;
  email: string;
  numeroEtudiant: string;
  statut: string;
  typeReclamation: string;
  objet: string;
  description: string;
  nomFichier: string;
}

const INIT: FormState = {
  prenom:"", nom:"", email:"", numeroEtudiant:"",
  statut:"", typeReclamation:"", objet:"",
  description:"", nomFichier:"",
};

export default function CentreReclamations() {
  const [form, setForm]       = useState<FormState>(INIT);
  const [loading, setLoading] = useState(false);
  const fileRef               = useRef<HTMLInputElement>(null);

  const set = (key: keyof FormState) =>
    (e: React.ChangeEvent
        HTMLInputElement | HTMLSelectElement | HTMLTextAreaElement>) =>
      setForm(f => ({ ...f, [key]: e.target.value }));

  const handleSubmit = async () => {
    const { prenom, nom, email, objet,
            description, statut, typeReclamation } = form;

    if (!prenom || !nom || !email || !objet ||
        !description || !statut || !typeReclamation) return;

    if (!/^[^\s@]+@[^\s@]+\.[^\s@]+$/.test(email)) return;

    setLoading(true);
    try {
      const id = await insertReclamation({
        prenom, nom, email,
        numeroEtudiant: form.numeroEtudiant || undefined,
        statut, typeReclamation, objet, description,
        nomFichier: form.nomFichier || undefined,
      });
      setForm(INIT);
      if (fileRef.current) fileRef.current.value = "";
    } finally {
      setLoading(false);
    }
  };

  return (
    <div>
      <h1>Centre de reclamations</h1>
      <button onClick={handleSubmit} disabled={loading}>
        {loading ? "Envoi en cours..." : "Envoyer la reclamation"}
      </button>
    </div>
  );
}
\end{lstlisting}

%===========================================
\newpage
\section*{Annexe C : Dashboard Monitor — Consultation des réclamations}
%===========================================

\subsection*{reclamation.mapper.ts}
\begin{lstlisting}[language=Java]
import type { ReclamationBackend, ReclamationUI }
    from "../model/reclamation.types";

const formatDate = (value: string | undefined): string => {
  if (!value) return "-";
  const parsed = new Date(value);
  if (Number.isNaN(parsed.getTime())) return value;
  return parsed.toLocaleDateString("fr-FR");
};

export const mapReclamationToUI = (
  backend: ReclamationBackend,
): ReclamationUI => ({
  id:              backend.id,
  nom:             `${backend.prenom} ${backend.nom}`.trim() || "-",
  email:           backend.email || "-",
  statut:          backend.statut || "-",
  typeReclamation: backend.typeReclamation || "-",
  objet:           backend.objet || "-",
  description:     backend.description || "-",
  date:            formatDate(backend.dateCreation),
});
\end{lstlisting}
